\chapter{Chapter06 — Examples}
%%% generated by the blueprint pipeline from FABL.Chapter06.Pseudorandomness.Examples
\section{Overview}

Declarations of the Lean module \texttt{FABL.Chapter06.Pseudorandomness.Examples}.

\section{Declarations}

\begin{theorem}[Fourier regularity of an affine subspace indicator]
\lean{FABL.isFourierRegular_setIndicator_binaryAffineSubspace}
\proofsource{Analysis-of-Boolean-Functions-by-Ryan-ODonnell}{
  pdf-ffbe9c7b3306-p144-b005,
  pdf-ffbe9c7b3306-p073-b002
}
\label{FABL.isFourierRegular_setIndicator_binaryAffineSubspace}
\leanok
\uses{FABL.F₂Cube, FABL.IsFourierRegular, FABL.abs_vectorFourierCoeff_setIndicator_binaryAffineSubspace_of_mem, FABL.abs_vectorFourierCoeff_setIndicator_binaryAffineSubspace_of_not_mem, FABL.binaryAffineSubspace, FABL.binaryFunctionOnSignCube, FABL.f₂CubeEquivFinset, FABL.f₂Support, FABL.inversePerpendicularCard, FABL.inversePerpendicularCard_pos, FABL.perpendicularSubspace, FABL.setIndicator, FABL.vectorFourierCoeff_eq_fourierCoeff_binaryFunctionOnSignCube, FABL.𝔽₂}
\difficulty{4}
Let $n\ge 0$, let $H$ be a linear subspace of $\mathbb F_2^n$, and let $a\in\mathbb
F_2^n$. The indicator of the affine subspace $a+H$, transported to the sign cube
$\{\pm1\}^n$ by the coordinatewise binary-to-sign correspondence, has every nonempty-set
Fourier coefficient bounded in absolute value by $2^{-d}$, where $d$ is the dimension of
the subspace orthogonal to $H$ under the standard dot product over $\mathbb F_2$.
\end{theorem}

\begin{theorem}[Affine subspace regularity by codimension]
\lean{FABL.isFourierRegular_setIndicator_binaryAffineSubspace_of_codimension}
\proofsource{Analysis-of-Boolean-Functions-by-Ryan-ODonnell}{
  pdf-ffbe9c7b3306-p073-b002,
  pdf-ffbe9c7b3306-p144-b005
}
\label{FABL.isFourierRegular_setIndicator_binaryAffineSubspace_of_codimension}
\leanok
\uses{FABL.F₂Cube, FABL.IsFourierRegular, FABL.binaryAffineSubspace, FABL.binaryFunctionOnSignCube, FABL.f₂Codimension, FABL.inversePerpendicularCard, FABL.isFourierRegular_setIndicator_binaryAffineSubspace, FABL.setIndicator, FABL.𝔽₂}
\difficulty{1}
Let $n,k\ge 0$, let $H$ be a linear subspace of $\mathbb F_2^n$, and let $a\in\mathbb
F_2^n$. Suppose the subspace orthogonal to $H$ under the standard dot product over
$\mathbb F_2$ has dimension $k$. Then the indicator of $a+H$, transported to
$\{\pm1\}^n$ by the coordinatewise binary-to-sign correspondence, has every nonempty-set
Fourier coefficient bounded in absolute value by $2^{-k}$.
\end{theorem}

\begin{theorem}[Regularity of the Mod-Two Inner-Product Indicator]
\lean{FABL.isFourierRegular_innerProductModTwo_zeroOne}
\proofsource{Analysis-of-Boolean-Functions-by-Ryan-ODonnell}{pdf-ffbe9c7b3306-p144-b005}
\label{FABL.isFourierRegular_innerProductModTwo_zeroOne}
\leanok
\uses{FABL.IsFourierRegular, FABL.abs_fourierCoeff_innerProductModTwoBoolean, FABL.binaryFunctionOnSignCube, FABL.binaryFunctionOnSignCube_booleanRealEmbedding_booleanFunctionF₂Encoding, FABL.booleanFunctionF₂Encoding, FABL.booleanRealEmbedding, FABL.fourierCoeff_one_sub_div_two, FABL.innerProductModTwoBoolean}
\difficulty{3}
For every natural number $m$, define a function on $\{\pm1\}^{2m}$ by decoding the
coordinates into two vectors $a,b\in\mathbb F_2^m$ and returning the indicator of
$\sum_{i=1}^m a_i b_i=1$ in $\mathbb F_2$. Every Fourier coefficient indexed by a
nonempty coordinate set has absolute value at most $2^{-(m+1)}$.
\end{theorem}

\begin{theorem}[Regularity of a Quadratic Boolean Function]
\lean{FABL.isFourierRegular_completeQuadratic_zeroOne}
\proofsource{Analysis-of-Boolean-Functions-by-Ryan-ODonnell}{pdf-ffbe9c7b3306-p144-b005}
\label{FABL.isFourierRegular_completeQuadratic_zeroOne}
\leanok
\uses{FABL.IsFourierRegular, FABL.abs_fourierCoeff_completeQuadraticBoolean, FABL.binaryFunctionOnSignCube, FABL.binaryFunctionOnSignCube_booleanRealEmbedding_booleanFunctionF₂Encoding, FABL.booleanFunctionF₂Encoding, FABL.booleanRealEmbedding, FABL.completeQuadraticBoolean, FABL.fourierCoeff_one_sub_div_two}
\difficulty{3}
Let $n$ be an even natural number, and define $Q(x)=\sum_{1\le i<j\le n}x_i x_j$ over
$\mathbb F_2$ for $x\in\mathbb F_2^n$. Regard the indicator of $Q(x)=1$ as a real-valued
function on $\{\pm1\}^n$ via the coordinatewise identification of bits with signs. Every
nonempty-set Fourier coefficient of this function has absolute value at most
$2^{-(n/2+1)}$.
\end{theorem}

\begin{theorem}[Parity is not Fourier regular below one]
\lean{FABL.parityFunction_not_isFourierRegular_of_lt_one}
\proofsource{Analysis-of-Boolean-Functions-by-Ryan-ODonnell}{pdf-ffbe9c7b3306-p144-b005}
\label{FABL.parityFunction_not_isFourierRegular_of_lt_one}
\leanok
\uses{FABL.BooleanFunction.toReal, FABL.IsFourierRegular, FABL.fourierCoeff_monomial, FABL.parityFunction, FABL.parityFunction_toReal}
\difficulty{3}
Let $S$ be a nonempty subset of the $n$ coordinates, and let $\varepsilon$ be a real
number with $\varepsilon<1$. The real-valued parity function on $S$ is not Fourier
regular with parameter $\varepsilon$.
\end{theorem}

\begin{theorem}[Fourier regularity of odd majority]
\lean{FABL.isFourierRegular_majority_odd}
\proofsource{Analysis-of-Boolean-Functions-by-Ryan-ODonnell}{pdf-ffbe9c7b3306-p144-b005}
\label{FABL.isFourierRegular_majority_odd}
\leanok
\uses{FABL.BooleanFunction.toReal, FABL.IsFourierRegular, FABL.abs_fourierCoeff_majority_le_singleton, FABL.fourierCoeff, FABL.influence, FABL.influence_eq_fourierCoeff_singleton_of_monotone, FABL.influence_le_one_div_sqrt_of_transitiveSymmetric_monotone, FABL.majority, FABL.majority_monotone, FABL.majority_transitiveSymmetric}
\difficulty{4}
For each integer $m\ge 0$, let $n=2m+1$ and consider the majority function on
$\{\pm1\}^n$, which outputs the sign of the sum of its coordinates (with nonnegative
sums assigned $+1$). Every nonempty-set Fourier coefficient of this function has
absolute value at most $1/\sqrt{n}$.
\end{theorem}

\begin{theorem}[Zero influences and constant functions]
\lean{FABL.hasSmallInfluences_zero_iff_exists_const}
\label{FABL.hasSmallInfluences_zero_iff_exists_const}
\leanok
\uses{FABL.HasSmallInfluences, FABL.const_hasSmallStableInfluences_zero_zero, FABL.fourierCoeff, FABL.hasSmallInfluences_iff, FABL.influence, FABL.influence_eq_sum_sq_fourierCoeff, FABL.influence_nonneg, FABL.isFourierRegular_zero_iff_exists_const}
\difficulty{4}
Let $n$ be a natural number and let $f$ be a real-valued function on the $n$-dimensional
sign cube. All influences of $f$ are at most zero if and only if $f$ is constant,
meaning that there exists $c\in\mathbb R$ such that $f(x)=c$ for every input $x$.
\end{theorem}

\begin{theorem}[Majority has small influences]
\lean{FABL.majority_hasSmallInfluences}
\statementsource{Analysis-of-Boolean-Functions-by-Ryan-ODonnell}{
  pdf-ffbe9c7b3306-p144-b005,
  pdf-ffbe9c7b3306-p062-b012
}
\label{FABL.majority_hasSmallInfluences}
\leanok
\uses{FABL.BooleanFunction.toReal, FABL.HasSmallInfluences, FABL.hasSmallInfluences_iff, FABL.influence_le_one_div_sqrt_of_transitiveSymmetric_monotone, FABL.majority, FABL.majority_monotone, FABL.majority_transitiveSymmetric}
\difficulty{1}
% dataset-pipeline: stub -- prose pending build_informal_statements.py
Statement pending regeneration.
\end{theorem}

\begin{theorem}[Stable influence of a parity function]
\lean{FABL.stableInfluence_parityFunction}
\proofsource{Analysis-of-Boolean-Functions-by-Ryan-ODonnell}{pdf-ffbe9c7b3306-p146-b001}
\label{FABL.stableInfluence_parityFunction}
\leanok
\uses{FABL.BooleanFunction.toReal, FABL.parityFunction, FABL.parityFunction_toReal, FABL.stableInfluence, FABL.stableInfluence_monomial}
\difficulty{1}
Let $n$ be a natural number, let $S$ be a subset of the coordinates, let $i$ be a
coordinate, and let $\rho\in\mathbb R$. The parity function on $S$ is the product of
the coordinates in $S$. Its stable influence at parameter $\rho$ and coordinate $i$
equals $\rho^{|S|-1}$ if $i\in S$, and equals $0$ otherwise.
\end{theorem}

\begin{theorem}[Stable influences of a parity function]
\lean{FABL.parityFunction_hasSmallStableInfluences_of_log_bound}
\proofsource{Analysis-of-Boolean-Functions-by-Ryan-ODonnell}{pdf-ffbe9c7b3306-p146-b001}
\label{FABL.parityFunction_hasSmallStableInfluences_of_log_bound}
\leanok
\uses{FABL.BooleanFunction.toReal, FABL.HasSmallStableInfluences, FABL.monomial_hasSmallStableInfluences_of_card, FABL.parityFunction, FABL.parityFunction_toReal}
\difficulty{5}
Let $n$ be a natural number, let $S$ be a subset of the $n$ coordinates, and let
$\varepsilon,\delta\in\mathbb R$ satisfy $\varepsilon>0$ and $0<\delta\le 1$. If
$|S|\ge \log(e/\varepsilon)/\delta$, then every coordinate has stable influence at
noise parameter $1-\delta$ at most $\varepsilon$ for the real-valued parity function
associated to $S$.
\end{theorem}

\begin{theorem}[Stable influence lower bound for balanced functions]
\lean{FABL.totalStableInfluence_one_sub_lower_of_dependsOn}
\label{FABL.totalStableInfluence_one_sub_lower_of_dependsOn}
\leanok
\uses{FABL.BooleanFunction, FABL.BooleanFunction.toReal, FABL.IsBalanced, FABL.dependsOn_toReal_iff, FABL.fourierCoeff, FABL.fourierCoeff_eq_zero_of_dependsOn_of_not_subset, FABL.stableInfluence, FABL.totalStableInfluence, FABL.totalStableInfluence_eq_sum_card_mul_rho_pow_mul_sq_fourierCoeff, FABL.variance, FABL.variance_eq_four_mul_probabilities, FABL.variance_eq_sum_sq_fourierCoeff}
\difficulty{6}
Let $n$ be a nonnegative integer, let $f:\{-1,1\}^n\to\{-1,1\}$ be a Boolean
function, and let $J\subseteq\{1,\ldots,n\}$. Suppose $f$ depends only on the
coordinates in $J$, and let $k$ be a natural number with $|J|\le k$. If $0<\delta<1$
and the real-valued sign representation of $f$ has mean zero, then
$(1-\delta)^{k-1}\le\sum_{i\in J}\operatorname{Inf}_{1-\delta}(f,i)$, where
$\operatorname{Inf}_{1-\delta}(f,i)$ denotes the stable influence at coordinate $i$.
\end{theorem}

\begin{theorem}[Stable influence of a balanced junta]
\lean{FABL.exists_stableInfluence_ge_of_isKJunta_of_balanced}
\proofsource{Analysis-of-Boolean-Functions-by-Ryan-ODonnell}{pdf-ffbe9c7b3306-p146-b002}
\label{FABL.exists_stableInfluence_ge_of_isKJunta_of_balanced}
\leanok
\uses{FABL.BooleanFunction, FABL.BooleanFunction.toReal, FABL.IsBalanced, FABL.IsKJunta, FABL.stableInfluence, FABL.totalStableInfluence_one_sub_lower_of_dependsOn}
\difficulty{4}
Let $f$ be a Boolean function on the $n$-dimensional sign cube, and let $k$ be a
positive integer. Suppose $f$ depends on at most $k$ coordinates and has mean zero. For
every $\delta\in(0,1)$, some coordinate $i$ has stable influence at parameter
$1-\delta$ at least $\frac{(1-\delta)^{k-1}}{k}$.
\end{theorem}

\begin{theorem}[Stable influences of balanced juntas]
\lean{FABL.not_hasSmallStableInfluences_of_isKJunta_of_balanced}
\proofsource{Analysis-of-Boolean-Functions-by-Ryan-ODonnell}{
  pdf-ffbe9c7b3306-p145-b008,
  pdf-ffbe9c7b3306-p146-b002
}
\label{FABL.not_hasSmallStableInfluences_of_isKJunta_of_balanced}
\leanok
\uses{FABL.BooleanFunction, FABL.BooleanFunction.toReal, FABL.HasSmallStableInfluences, FABL.IsBalanced, FABL.IsKJunta, FABL.exists_stableInfluence_ge_of_isKJunta_of_balanced}
\difficulty{4}
Let $f$ be a sign-valued Boolean function on the $n$-dimensional cube, and suppose it
depends on at most $k$ coordinates, where $k>0$. If the mean of its real-valued sign
representation is zero and $0<\delta<1$, then it is not the case that every coordinate
has stable influence at parameter $1-\delta$ at most $(1-\delta)^k/k$.
\end{theorem}

\begin{definition}[Leading-coordinate sign extension]
\lean{FABL.leadingCoordinateTimes}
\label{FABL.leadingCoordinateTimes}
\leanok
\uses{FABL.BooleanFunction}
Let $g$ be a function from sign vectors of length $n$ to signs. Define a function on
sign vectors of length $n+1$ by multiplying the first coordinate by the value of $g$ on
the vector formed from the remaining coordinates.
\end{definition}

\begin{theorem}[Leading-coordinate product on a prefixed input]
\lean{FABL.leadingCoordinateTimes_fin_cons}
\label{FABL.leadingCoordinateTimes_fin_cons}
\leanok
\uses{FABL.BooleanFunction, FABL.Sign, FABL.leadingCoordinateTimes}
\difficulty{1}
Let $g:\{−1,1\}^n\to\{−1,1\}$, let $b\in\{−1,1\}$, and let $x\in\{−1,1\}^n$. The
function on the $(n+1)$-dimensional sign cube formed by multiplying the first coordinate
by $g$ of the remaining coordinates, evaluated at the input with first coordinate $b$
and remaining coordinates $x$, equals $b g(x)$.
\end{theorem}

\begin{theorem}[Real-valued leading-coordinate product]
\lean{FABL.leadingCoordinateTimes_toReal_fin_cons}
\label{FABL.leadingCoordinateTimes_toReal_fin_cons}
\leanok
\uses{FABL.BooleanFunction, FABL.BooleanFunction.toReal, FABL.Sign, FABL.leadingCoordinateTimes, FABL.leadingCoordinateTimes_fin_cons, FABL.signValue}
\difficulty{3}
Let $g:\{−1,1\}^n\to\{−1,1\}$, let $b\in\{−1,1\}$, and let $x\in\{−1,1\}^n$. The
real-valued version of the leading-coordinate product function at the input $(b,x)$
equals the real value of $b$ multiplied by the real-valued version of $g$ at $x$.
\end{theorem}

\begin{theorem}[Fourier coefficient of a leading-coordinate product]
\lean{FABL.fourierCoeff_leadingCoordinateTimes_insert_zero_tailFrequency}
\label{FABL.fourierCoeff_leadingCoordinateTimes_insert_zero_tailFrequency}
\leanok
\uses{FABL.BooleanFunction, FABL.BooleanFunction.toReal, FABL.firstCoordinateSlice, FABL.firstCoordinateSlice_apply, FABL.fourierCoeff, FABL.fourierCoeff_insert_zero_tailFrequency, FABL.leadingCoordinateTimes, FABL.leadingCoordinateTimes_toReal_fin_cons, FABL.monomial, FABL.tailFrequency}
\difficulty{4}
Let $n$ be a nonnegative integer, let $g:\{\pm1\}^n\to\{\pm1\}$, and let $S$ be a subset
of the $n$ coordinates. Define a function on $n+1$ coordinates by multiplying the first
coordinate by $g$ evaluated on the remaining coordinates. Its Fourier coefficient at the
frequency consisting of coordinate $0$ together with the coordinates corresponding to
$S$ equals the Fourier coefficient of $g$ at $S$.
\end{theorem}

\begin{theorem}[Stability bounded by a stable influence]
\lean{FABL.stabilityCurve_le_stableInfluence_leadingCoordinateTimes}
\label{FABL.stabilityCurve_le_stableInfluence_leadingCoordinateTimes}
\leanok
\uses{FABL.BooleanFunction, FABL.BooleanFunction.toReal, FABL.card_tailFrequency, FABL.fourierCoeff, FABL.fourierCoeff_leadingCoordinateTimes_insert_zero_tailFrequency, FABL.leadingCoordinateTimes, FABL.stabilityCurve, FABL.stableInfluence, FABL.tailFrequency, FABL.zero_notMem_tailFrequency}
\difficulty{5}
Let $g$ be a Boolean function on the $n$-dimensional sign cube and let $\rho\in[0,1]$.
Form a Boolean function on the $(n+1)$-dimensional sign cube by multiplying the first
coordinate by $g$ evaluated on the remaining coordinates. The stability curve of the
real-valued version of $g$ at $\rho$ is at most the stable influence of this new
function in its first coordinate, where stable influence weights each squared Fourier
coefficient containing that coordinate by $\rho$ raised to one less than the size of its
index set.
\end{theorem}

\begin{theorem}[An arcsine bound]
\lean{FABL.one_sub_sqrt_le_two_div_pi_mul_arcsin_one_sub}
\label{FABL.one_sub_sqrt_le_two_div_pi_mul_arcsin_one_sub}
\leanok
\uses{FABL.two_div_pi_mul_arcsin_eq_one_sub_two_div_pi_mul_arccos}
\difficulty{6}
For every real number $\delta$ with $0\leq\delta\leq 1$, one has $1-\sqrt{\delta}\leq
\frac{2}{\pi}\arcsin(1-\delta)$.
\end{theorem}

\begin{theorem}[Stable influence of leading-coordinate majority]
\lean{FABL.one_sub_sqrt_le_stableInfluence_leadingCoordinateTimes_majority_odd}
\label{FABL.one_sub_sqrt_le_stableInfluence_leadingCoordinateTimes_majority_odd}
\leanok
\uses{FABL.BooleanFunction.toReal, FABL.leadingCoordinateTimes, FABL.majority, FABL.noiseStability, FABL.one_sub_sqrt_le_two_div_pi_mul_arcsin_one_sub, FABL.stabilityCurve, FABL.stabilityCurve_eq_noiseStability, FABL.stabilityCurve_le_stableInfluence_leadingCoordinateTimes, FABL.stableInfluence, FABL.sum_sq_fourierCoeff_eq_one, FABL.two_div_pi_mul_arcsin_le_noiseStability_majority_odd}
\difficulty{5}
Let $m$ be a natural number and let $0\leq\delta\leq1$. The stable influence at
parameter $1-\delta$ of the leading coordinate of the real-valued function obtained by
multiplying odd majority on $2m+1$ signs by a new leading coordinate is at least
$1-\sqrt{\delta}$.
\end{theorem}

\begin{definition}[Parity-weighted Boolean function]
\lean{FABL.parityTimes}
\statementsource{Analysis-of-Boolean-Functions-by-Ryan-ODonnell}{
  pdf-ffbe9c7b3306-p023-b001,
  pdf-ffbe9c7b3306-p147-b006,
  pdf-ffbe9c7b3306-p147-b007
}
\label{FABL.parityTimes}
\leanok
\uses{FABL.BooleanFunction, FABL.parityFunction}
For a Boolean function $g$ on the $n$-dimensional sign cube and a set of coordinates
$S\subseteq\{0,\ldots,n-1\}$, define the function whose value at each input $x$ is the
parity function on $S$ evaluated at $x$, multiplied by $g(x)$.
\end{definition}

\begin{theorem}[Real form of parity multiplication]
\lean{FABL.parityTimes_toReal}
\label{FABL.parityTimes_toReal}
\leanok
\uses{FABL.BooleanFunction, FABL.BooleanFunction.toReal, FABL.monomial, FABL.parityFunction, FABL.parityFunction_toReal, FABL.parityTimes, FABL.signValue}
\difficulty{3}
Let $n$ be a natural number, let $S$ be a set of coordinates among $n$ coordinates, and
let $g$ be a Boolean function on the $n$-dimensional sign cube. The real-valued function
associated with the pointwise product of $g$ and the parity on $S$ is the pointwise
product of the monomial on $S$ and the real-valued function associated with $g$.
\end{theorem}

\begin{theorem}[Fourier coefficients after parity multiplication]
\lean{FABL.fourierCoeff_parityTimes}
\label{FABL.fourierCoeff_parityTimes}
\leanok
\uses{FABL.BooleanFunction, FABL.BooleanFunction.toReal, FABL.fourierCoeff, FABL.monomial_mul_monomial, FABL.parityTimes, FABL.parityTimes_toReal}
\difficulty{3}
Let $n$ be a nonnegative integer, let $S,T\subseteq\{1,\ldots,n\}$, and let $g$ be a
function from the sign cube to $\{\pm1\}$. Multiply $g$ pointwise by the parity
character $\prod_{i\in S}x_i$. The Fourier coefficient of the resulting function at
$T$ equals the Fourier coefficient of $g$ at the symmetric difference $S\triangle T$.
\end{theorem}

\begin{theorem}[Resilience after multiplying by parity]
\lean{FABL.parityTimes_isResilient_of_dependsOn_compl}
\proofsource{Analysis-of-Boolean-Functions-by-Ryan-ODonnell}{
  pdf-ffbe9c7b3306-p147-b007,
  pdf-ffbe9c7b3306-p147-b006
}
\label{FABL.parityTimes_isResilient_of_dependsOn_compl}
\leanok
\uses{FABL.BooleanFunction, FABL.BooleanFunction.toReal, FABL.IsResilient, FABL.dependsOn_toReal_iff, FABL.fourierCoeff, FABL.fourierCoeff_eq_zero_of_dependsOn_of_not_subset, FABL.fourierCoeff_parityTimes, FABL.isBalanced_iff_fourierCoeff_empty_eq_zero, FABL.parityTimes}
\difficulty{5}
Let $n,k$ be natural numbers, let $S$ be a set of coordinates in $\{0,\ldots,n-1\}$ with
$|S|=k+1$, and let $g:\{−1,1\}^n\to\{−1,1\}$. Suppose $g$ depends only on coordinates
outside $S$. Then the function obtained by multiplying $g$ pointwise by the parity of
the coordinates in $S$ is $k$-resilient: it is balanced and all its nonconstant Fourier
coefficients indexed by sets of size at most $k$ vanish.
\end{theorem}

\begin{definition}[First two thirds of the coordinates]
\lean{FABL.firstTwoThirds}
\statementsource{Analysis-of-Boolean-Functions-by-Ryan-ODonnell}{pdf-ffbe9c7b3306-p147-b007}
\label{FABL.firstTwoThirds}
\leanok
For a positive integer $m$, define a subset of the coordinates of a $3m$-dimensional
cube to consist of the indices $i$ with $0\le i<2m$.
\end{definition}

\begin{definition}[Last two thirds of the coordinates]
\lean{FABL.lastTwoThirds}
\statementsource{Analysis-of-Boolean-Functions-by-Ryan-ODonnell}{pdf-ffbe9c7b3306-p147-b007}
\label{FABL.lastTwoThirds}
\leanok
For a positive integer $m$, define a subset of the coordinates of a $3m$-dimensional
cube to consist of the indices $i$ with $m\le i<3m$.
\end{definition}

\begin{theorem}[Size of the first two-thirds set]
\lean{FABL.card_firstTwoThirds}
\label{FABL.card_firstTwoThirds}
\leanok
\uses{FABL.firstTwoThirds}
\difficulty{1}
For every positive natural number $m$, the set of the first $2m$ indices among $3m$
indices has cardinality $2m$.
\end{theorem}

\begin{theorem}[Cardinality of the last two thirds]
\lean{FABL.card_lastTwoThirds}
\label{FABL.card_lastTwoThirds}
\leanok
\uses{FABL.lastTwoThirds}
\difficulty{2}
For a positive integer $m$, the set of integers in $\{0,\ldots,3m-1\}$ that are at least
$m$ has cardinality $2m$.
\end{theorem}

\begin{theorem}[Symmetric difference of two index sets]
\lean{FABL.card_firstTwoThirds_symmDiff_lastTwoThirds}
\label{FABL.card_firstTwoThirds_symmDiff_lastTwoThirds}
\leanok
\uses{FABL.firstTwoThirds, FABL.lastTwoThirds}
\difficulty{5}
For every positive natural number $m$, the symmetric difference of the first $2m$
indices and the last $2m$ indices among $3m$ indices has cardinality $2m$.
\end{theorem}

\begin{definition}[A Boolean function from two parities]
\lean{FABL.correlationImmuneAndExample}
\statementsource{Analysis-of-Boolean-Functions-by-Ryan-ODonnell}{
  pdf-ffbe9c7b3306-p147-b006,
  pdf-ffbe9c7b3306-p147-b007,
  pdf-ffbe9c7b3306-p170-b002
}
\label{FABL.correlationImmuneAndExample}
\leanok
\uses{FABL.BooleanFunction, FABL.andFunction, FABL.firstTwoThirds, FABL.lastTwoThirds, FABL.parityFunction}
For a positive integer $m$, define a Boolean function on the $3m$-dimensional sign cube
by applying the two-input AND function to the parities on the first $2m$ coordinates and
on the last $2m$ coordinates.
\end{definition}

\begin{theorem}[Two-variable AND formula]
\lean{FABL.andFunction_two_toReal}
\label{FABL.andFunction_two_toReal}
\leanok
\uses{FABL.Sign, FABL.andFunction, FABL.signValue}
\difficulty{4}
For signs $a,b\in\{-1,1\}$, the real value of the two-variable AND function at $(a,b)$
is $\frac{1+a+b-ab}{2}$.
\end{theorem}

\begin{theorem}[Real-valued formula for the parity-based function]
\lean{FABL.correlationImmuneAndExample_toReal}
\label{FABL.correlationImmuneAndExample_toReal}
\leanok
\uses{FABL.BooleanFunction.toReal, FABL.andFunction_two_toReal, FABL.correlationImmuneAndExample, FABL.firstTwoThirds, FABL.lastTwoThirds, FABL.monomial, FABL.monomial_mul_monomial, FABL.parityFunction, FABL.parityFunction_toReal, FABL.signValue}
\difficulty{3}
For every positive integer $m$, the real-valued version of the function obtained by
taking the AND of the parities on the first and last $2m$ coordinates is, at each input
$x\in\{\!-1,1\!\}^{3m}$, equal to one half of $1$ plus the monomial on the first
coordinate set plus the monomial on the last coordinate set minus the monomial on their
symmetric difference.
\end{theorem}

\begin{theorem}[Fourier coeff correlation immune and example]
\lean{FABL.fourierCoeff_correlationImmuneAndExample}
\statementsource{Analysis-of-Boolean-Functions-by-Ryan-ODonnell}{pdf-ffbe9c7b3306-p147-b007}
\label{FABL.fourierCoeff_correlationImmuneAndExample}
\leanok
\uses{FABL.BooleanFunction.toReal, FABL.correlationImmuneAndExample, FABL.correlationImmuneAndExample_toReal, FABL.expect_monomial, FABL.expect_monomial_mul, FABL.firstTwoThirds, FABL.fourierCoeff, FABL.lastTwoThirds, FABL.monomial}
\difficulty{4}
% dataset-pipeline: stub -- prose pending build_informal_statements.py
Statement pending regeneration.
\end{theorem}

\begin{theorem}[Correlation immunity of a parity-based function]
\lean{FABL.correlationImmuneAndExample_isCorrelationImmune}
\proofsource{Analysis-of-Boolean-Functions-by-Ryan-ODonnell}{
  pdf-ffbe9c7b3306-p147-b006,
  pdf-ffbe9c7b3306-p147-b007,
  pdf-ffbe9c7b3306-p170-b002
}
\label{FABL.correlationImmuneAndExample_isCorrelationImmune}
\leanok
\uses{FABL.IsCorrelationImmune, FABL.card_firstTwoThirds_symmDiff_lastTwoThirds, FABL.correlationImmuneAndExample, FABL.firstTwoThirds, FABL.fourierCoeff_correlationImmuneAndExample, FABL.lastTwoThirds}
\difficulty{5}
For a positive integer $m$, define a function on $\{\!-1,1\!\}^{3m}$ by taking the AND
of the parities on the first $2m$ coordinates and on the last $2m$ coordinates. Every
nonempty coordinate set of size at most $2m-1$ has Fourier coefficient zero for this
function; hence it is correlation immune of order $2m-1$.
\end{theorem}

\begin{theorem}[Mean correlation immune and example]
\lean{FABL.mean_correlationImmuneAndExample}
\statementsource{Analysis-of-Boolean-Functions-by-Ryan-ODonnell}{pdf-ffbe9c7b3306-p147-b007}
\label{FABL.mean_correlationImmuneAndExample}
\leanok
\uses{FABL.BooleanFunction.toReal, FABL.card_firstTwoThirds_symmDiff_lastTwoThirds, FABL.correlationImmuneAndExample, FABL.firstTwoThirds, FABL.fourierCoeff_correlationImmuneAndExample, FABL.lastTwoThirds, FABL.mean, FABL.mean_eq_fourierCoeff_empty}
\difficulty{4}
% dataset-pipeline: stub -- prose pending build_informal_statements.py
Statement pending regeneration.
\end{theorem}

\begin{theorem}[Failure of resilience for the parity-based function]
\lean{FABL.correlationImmuneAndExample_not_isResilient}
\proofsource{Analysis-of-Boolean-Functions-by-Ryan-ODonnell}{
  pdf-ffbe9c7b3306-p147-b007,
  pdf-ffbe9c7b3306-p147-b006
}
\label{FABL.correlationImmuneAndExample_not_isResilient}
\leanok
\uses{FABL.IsBalanced, FABL.IsResilient, FABL.correlationImmuneAndExample, FABL.mean_correlationImmuneAndExample}
\difficulty{3}
For a positive integer $m$, define a function on $\{\!-1,1\!\}^{3m}$ by taking the AND
of the parities on the first $2m$ coordinates and on the last $2m$ coordinates. This
function is not resilient of order $2m-1$, meaning that it does not both have all
nonempty Fourier coefficients of degree at most $2m-1$ equal to zero and have mean zero.
\end{theorem}

\begin{theorem}[Uniform probability correlation immune and example eq true]
\lean{FABL.uniformProbability_correlationImmuneAndExample_eq_true}
\statementsource{Analysis-of-Boolean-Functions-by-Ryan-ODonnell}{pdf-ffbe9c7b3306-p147-b007}
\label{FABL.uniformProbability_correlationImmuneAndExample_eq_true}
\leanok
\uses{FABL.correlationImmuneAndExample, FABL.mean_correlationImmuneAndExample, FABL.mean_eq_probability_one_sub_probability_neg_one, FABL.uniformProbability, FABL.uniformProbability_one_add_neg_one_eq_one}
\difficulty{3}
% dataset-pipeline: stub -- prose pending build_informal_statements.py
Statement pending regeneration.
\end{theorem}
